\begin{center}
    {\LARGE \textbf{2012无机化学}} \\[2ex]
    {\large 时量180分钟 \quad 满分100分}
\end{center}

\vspace{0.5cm}
\noindent\textbf{注意事项:}
\begin{enumerate}[label=\arabic*., parsep=0pt, itemsep=0pt]
    \item 答题前,考生先将自己的姓名、学号、考场号、座位号填写在答题卡上,将条形码准确粘贴在考生信息条形码粘贴区。
    \item 考试结束后,将本试卷与答题纸一并收回。
    \item 回答各个问题时,将答案写在答题纸上,写在本试卷上无效。
\end{enumerate}

\vspace{0.5cm}
\noindent\textbf{一、完成并配平下列反应的方程式（24 pts.）}

\begin{enumerate}[label=\arabic*., start=1, itemsep=1.5ex]
    \item 向酸性的\ce{KI}溶液中通入\ce{O3}。
    \item 将\ce{PbO2}溶于\ce{HCl}溶液中。
    \item 向\ce{CuSO4}溶液中滴加\ce{KCN}溶液。
    \item 向次磷酸溶液中滴加硝酸银溶液。
    \item \ce{MnSO4}溶液与\ce{NaBiO3}混合后加入稀硫酸。
    \item 向\ce{NiCl2}溶液中加入\ce{NaOH}溶液和溴水。
\end{enumerate}

\vspace{0.5cm}
\noindent\textbf{二、解释下列实验事实（16 pts.）}

\begin{enumerate}[label=\arabic*., start=1, itemsep=1.5ex]
    \item 制备少量\ce{NO2}时，使用\ce{Pb(NO3)2}热分解而不使用\ce{NaNO3}。
    \item 向\ce{FeSO4}溶液中加入碘水，碘水不褪色；再加入\ce{NH4F}溶液，则碘水褪色。
    \item 将\ce{SO2}通入稀的品红溶液，溶液的红色消失；过一段时间后，溶液的红色又逐渐恢复。
    \item 向\ce{NiSO4}溶液中滴加\ce{NaOH}溶液生成绿色沉淀，再加入\ce{H2O2}溶液沉淀不变色；但用氯水代替\ce{H2O2}溶液则沉淀变色。
\end{enumerate}

\vspace{0.5cm}
\noindent\textbf{三、回答下列各题（54 pts.）}

\begin{enumerate}[label=\arabic*., start=1, itemsep=1.5ex]
    \item 写出元素\ce{Cr}、\ce{Pt}的电子结构式，并用四个量子数表示出这两种元素价层所有电子的运动状态。
    \item 已知
    \[
    \begin{array}{lccc}
    & \ce{NaCl} & \ce{MgO} & \ce{Al2O3} \\
    \text{熔点}/^\circ\text{C} & 801 & 2806 & 2054
    \end{array}
    \]
    试说明为什么\ce{MgO}的熔点既比\ce{NaCl}的高，又比\ce{Al2O3}的高？
    \item 比较下列各组离子或化合物的稳定性，并说明原因。
    \begin{enumerate}[label=(\arabic*), itemsep=1ex]
        \item \ce{NaHCO3}与\ce{Na2CO3}
        \item \ce{CaSO4}与\ce{ZnSO4}
        \item \ce{[Fe(CN)6]^{4-}}与\ce{[Fe(CN)6]^{3-}}
        \item \ce{[CuCl3]^{-}}与\ce{[Cu(NH3)4]^{2+}}
    \end{enumerate}
    \item 为什么\ce{NH3}对于氢离子的配位能力比\ce{PH3}强，而\ce{PH3}对于过渡金属离子的配位能力比\ce{NH3}强？
    \item 试写出\ce{VA}族元素三氯化物的水解反应方程式，并简要总结反应规律。
    \item 为什么\ce{[CoCl4]^{2-}}和\ce{[NiCl4]^{2-}}为四面体结构，而\ce{[CuCl4]^{2-}}和\ce{[PtCl4]^{2-}}却为正方形结构？
    \item 命名下列配位化合物，指出中心离子的杂化类型，并画出配位单元可能存在的几何异构体。
    \begin{enumerate}[label=(\arabic*), itemsep=1ex]
        \item \ce{[Pt(NH3)2Cl2(OH)2]}
        \item \ce{[Co(NH3)(en)Cl3]}
    \end{enumerate}
    \item 通过计算说明\ce{Cu+}在氨水中能否稳定存在？
    已知：
    \[
    \begin{array}{ll}
    E^\ominus(\ce{Cu^{2+}/Cu}) = 0.34\text{V}, & E^\ominus(\ce{Cu+/Cu}) = 0.52\text{V}, \\
    \lg K^\ominus_{\text{稳}[\ce{Cu(NH3)4}^{2+}]} = 13.32, & \lg K^\ominus_{\text{稳}[\ce{Cu(NH3)2+}]} = 10.86
    \end{array}
    \]
\end{enumerate}

\vspace{0.5cm}
\noindent\textbf{四、制备与分离（26 pts.）}

\begin{enumerate}[label=\arabic*., start=1, itemsep=1.5ex]
    \item 用方程式表示以\ce{NH4HSO4}为原料，利用电解-水解法制备\ce{H2O2}的过程。
    \item 工业上使用\ce{KCl}为原料制备\ce{KClO3}。写出相关的反应方程式。
    \item 工业碘中常含有\ce{ICl}或\ce{IBr}，如何分离提纯？写出相关反应方程式。
    \item 设计方案分离含有\ce{Ag+}、\ce{Cu^{2+}}、\ce{Zn^{2+}}、\ce{Cd^{2+}}、\ce{Hg^{2+}}的混合溶液。
\end{enumerate}

\vspace{0.5cm}
\noindent\textbf{五、推断（30 pts.）}

\begin{enumerate}[label=\arabic*., start=1, itemsep=1.5ex]
    \item 灰黑色化合物\ce{A}溶于浓盐酸形成粉红色溶液，放出气体\ce{B}，从溶液中析出粉红色物质\ce{C}，\ce{C}受热脱水转化为蓝色物质\ce{D}。将\ce{C}的浓溶液用醋酸酸化并加入亚硝酸钾，则析出黄色沉淀\ce{E}。在\ce{C}的溶液中加入氯化铵、浓氨水、过氧化氢，加热反应，冷却后析出橙黄色晶体\ce{F}。\ce{D}加到氨水中，生成粉红色沉淀\ce{G}，当通入气体\ce{B}时，\ce{G}转化为黑色沉淀\ce{H}。\ce{H}在盐酸中与氧化亚锡作用生成\ce{C}的粉红色溶液。
    试给出\ce{A}、\ce{B}、\ce{C}、\ce{D}、\ce{E}、\ce{F}、\ce{G}和\ce{H}的化学式，并完成各步的化学反应方程式。
    \item 白色固体\ce{A}与水混合后生成白色沉淀\ce{B}。\ce{B}溶于\ce{HCl}中得无色溶液\ce{C}。向\ce{C}中加入\ce{NaOH}溶液有白色沉淀\ce{D}生成，\ce{NaOH}溶液过量则\ce{D}溶解得到无色溶液\ce{E}。\ce{B}溶于\ce{HCl}后缓慢滴加到\ce{HgCl2}溶液中先有白色沉淀\ce{E}生成，而后白色沉淀逐渐变灰，最后转化为黑色沉淀\ce{G}。\ce{B}溶于稀\ce{HNO3}后加入\ce{AgNO3}溶液得到白色沉淀\ce{H}。\ce{H}溶于氨水得无色溶液，加\ce{HNO3}酸化又产生白色沉淀\ce{H}。
    试给出\ce{A}、\ce{B}、\ce{C}、\ce{D}、\ce{E}、\ce{F}、\ce{G}和\ce{H}所代表的物质的化学式。写出相关的反应方程式。
    \item 化合物\ce{A}是一种黑色固体，它不溶于水，也不溶于稀\ce{HAc}及稀碱溶液，而易溶于热\ce{HCl}中，生成一种绿色的溶液\ce{B}；如溶液\ce{B}与铜丝一起煮，即逐渐变为土黄色溶液\ce{C}；溶液\ce{C}若用大量水稀释时会生成白色沉淀\ce{D}；\ce{D}可溶于氨溶液中生成无色溶液\ce{E}；\ce{E}暴露于空气中则迅速变成深蓝色溶液\ce{F}；往\ce{F}中加入过量\ce{KCN}时蓝色消失，生成溶液\ce{G}；往\ce{G}中加入锌粒，则生成红色沉淀\ce{H}；\ce{H}不溶于酸和稀碱中，但可溶于稀\ce{HNO3}中生成浅蓝色溶液\ce{I}；往\ce{I}中慢慢加入\ce{NaOH}溶液则生成天蓝色沉淀\ce{J}，如将\ce{J}过滤取出后强热又生成原来的化合物\ce{A}。
    问\ce{A}、\ce{B}、\ce{C}、\ce{D}、\ce{E}、\ce{F}、\ce{G}、\ce{H}、\ce{I}、\ce{J}各为何物？
\end{enumerate}